% symbols.tex
% Appendix -- Symbols and meanings.  Reference list of mathematical
% symbols introduced in this paper.  Update this list whenever a new
% standing symbol enters the body of the paper.
%
% Section column is the section where the symbol is first used in its
% standing meaning.  Many symbols are reused in later sections; the
% Section column points to the *definition*, not the most recent use.

\label{sec:symbols}

\placeholder{Update entries whenever a new standing symbol enters the
body of the paper.  Many entries below are placeholders pinned to
their planned section; the meaning is fixed by the
\texttt{scope-and-outline} note (\texttt{notes/dcl\_mathematics\_scope\_and\_outline.md}).}

{\small
\setlength{\tabcolsep}{6pt}
\renewcommand{\arraystretch}{1.10}
\begin{longtable}{@{}p{2.6cm}p{8.4cm}p{3.0cm}@{}}
\caption{\small Symbols and their meanings, with the section where
each symbol is first defined.  Symbols introduced only inside a
single proof or example (e.g.\ dummy summation indices) are not
listed.}
\label{tab:symbols} \\
\hline
\textbf{Symbol} & \textbf{Meaning} & \textbf{Section} \\
\hline
\endfirsthead

\multicolumn{3}{l}{\textit{(Table~\ref{tab:symbols}, continued from previous page)}} \\
\hline
\textbf{Symbol} & \textbf{Meaning} & \textbf{Section} \\
\hline
\endhead

\hline
\multicolumn{3}{r}{\textit{(continued on next page)}} \\
\endfoot

\hline
\endlastfoot

% ── Lattice and ambient spaces ─────────────────────────────────────
\multicolumn{3}{@{}l}{\textit{Lattice and ambient spaces}} \\
$n$
    & Lattice dimension; free parameter throughout the paper.  Specialised to $n = 3$ only when matching Paper~I's physical setting.
    & \S\ref{sec:lattice_n_dim} \\
$\Lambda_n$
    & The bipartite octahedral lattice in dimension $n$.
    & \S\ref{sec:lattice_n_dim} \\
$V_n$
    & Vertex set of $\Lambda_n$, identified with $\mathbb{Z}^n$ (provisional).
    & \S\ref{subsec:lattice_definition} \\
$V_n^{+}, V_n^{-}$
    & The two colour classes of the bipartition; $V_n = V_n^{+} \sqcup V_n^{-}$.
    & \S\ref{subsec:lattice_parity} \\
$E_n$
    & Edge set of $\Lambda_n$: axial $\pm e_i$ edges plus octahedral cross-edges.
    & \S\ref{subsec:lattice_definition} \\
$\epsilon$
    & Lattice-spacing parameter ($\epsilon > 0$); $\epsilon \to 0$ gives the discrete-to-continuum limit.
    & \S\ref{subsec:limit_parameter} \\
$\epsilon \Lambda_n$
    & Rescaled lattice $\subset \mathbb{R}^n$ with spacing $\epsilon$.
    & \S\ref{subsec:limit_parameter} \\
$\mathbb{R}^n$, $\mathbb{Z}^n$
    & Continuum ambient space; integer lattice (used as vertex-set model for $\Lambda_n$).
    & \S\ref{subsec:lattice_definition} \\

% ── Symmetries ─────────────────────────────────────────────────────
\multicolumn{3}{@{}l}{\textit{Symmetries}} \\
$B_n$
    & Hyperoctahedral group: signed permutations of the $n$ coordinate axes; acts on $\Lambda_n$.
    & \S\ref{subsec:autom_generators} \\
$\mathfrak{a}_n$
    & The per-site automorphism algebra of $\Lambda_n$; provisional notation. At $n = 3$, $\dim \mathfrak{a}_3 = 71$ (catalogued in \texttt{dcl-zoo}).
    & \S\ref{sec:automorphism_algebra} \\

% ── Discrete differential geometry ─────────────────────────────────
\multicolumn{3}{@{}l}{\textit{Discrete differential geometry}} \\
$\Omega^k(\Lambda_n)$
    & Discrete $k$-forms on $\Lambda_n$: functions on oriented $k$-cells.
    & \S\ref{subsec:dec_forms} \\
$d$
    & Discrete exterior derivative $\Omega^k \to \Omega^{k+1}$.  Satisfies $d \circ d = 0$.
    & \S\ref{subsec:dec_d} \\
$d^\ast$
    & Discrete codifferential: formal adjoint of $d$ with respect to the inner product on forms.
    & \S\ref{subsec:dec_codifferential} \\
$\langle \cdot, \cdot \rangle$
    & Inner product on $\Omega^k(\Lambda_n)$ defined via the discrete measure.
    & \S\ref{subsec:dec_codifferential} \\
$\Delta = -d^\ast d - d d^\ast$
    & Discrete Hodge Laplacian on $\Omega^k(\Lambda_n)$.  Converges to the continuum Laplacian in the $\epsilon \to 0$ limit.
    & \S\ref{subsec:consistency_conditions} \\

% ── Function spaces ────────────────────────────────────────────────
\multicolumn{3}{@{}l}{\textit{Function spaces}} \\
$\ell^p(\Lambda_n)$
    & Discrete Lebesgue space, $1 \leq p \leq \infty$.
    & \S\ref{subsec:lp_spaces} \\
$\ell^{k,p}(\Lambda_n)$
    & Discrete Sobolev-analogue space: lattice functions with discrete derivatives (via $d$) up to order $k$ in $\ell^p$.
    & \S\ref{subsec:sobolev_analogue} \\
$\|f\|_{p, \epsilon}$
    & $\epsilon$-rescaled $\ell^p$-norm: $\|f\|_{p, \epsilon}^p = \epsilon^n \sum_{x} |f(x)|^p$.  Converges to $\|f\|_{L^p(\mathbb{R}^n)}$ in the $\epsilon \to 0$ limit.
    & \S\ref{subsec:embeddings} \\
$L^p(\mathbb{R}^n)$, $W^{k,p}(\mathbb{R}^n)$
    & Continuum Lebesgue and Sobolev spaces; targets of the discrete-to-continuum limit.
    & \S\ref{subsec:function_space_convergence} \\

% ── Measure theory ─────────────────────────────────────────────────
\multicolumn{3}{@{}l}{\textit{Measure theory}} \\
$\mu_\#$
    & Counting measure on $\Lambda_n$: $\mu_\#(A) = |A|$ for finite $A \subset \Lambda_n$.
    & \S\ref{subsec:counting_measure} \\
$\mu_{\#, \epsilon}$
    & $\epsilon$-rescaled counting measure on $\epsilon \Lambda_n$: $\mu_{\#, \epsilon} = \epsilon^n \mu_\#$.  Converges to Lebesgue measure on $\mathbb{R}^n$ in the limit.
    & \S\ref{subsec:counting_measure} \\

% ── Kinetic theory (motivating examples) ───────────────────────────
\multicolumn{3}{@{}l}{\textit{Kinetic theory (motivating examples, Section~\ref{sec:motivating_examples})}} \\
$V$
    & Discrete velocity space attached to each lattice site.  Provisionally, the neighbour-vector set of $\Lambda_n$.
    & \S\ref{subsec:boltzmann} \\
$f(x, v, t)$
    & Lattice kinetic distribution function: density of particles at $x \in \Lambda_n$ with velocity $v \in V$ at time $t$.
    & \S\ref{subsec:boltzmann} \\
$f_{\text{eq}}$
    & Local equilibrium (Maxwellian-analogue) distribution; parametrised by the conserved moments.
    & \S\ref{subsec:boltzmann} \\
$Q[f]$
    & Boltzmann-shaped collision operator on the lattice.
    & \S\ref{subsec:boltzmann} \\
$H[f]$
    & Lattice $H$-functional: $H[f] = \sum_{x, v} f(x, v) \log f(x, v)$.  Conjectured non-increasing under the lattice kinetic dynamics.
    & \S\ref{subsec:boltzmann} \\
$\rho$, $u$, $e$
    & Fluid density, velocity, energy (the conserved moments of $f$); arguments of $f_{\text{eq}}$.
    & \S\ref{subsec:fluid_limit} \\

\hline
\end{longtable}
}
